% Slot 1: C1--C4 (Thu 26 Nov, 09:00--11:00 Rome)

\section{Slot 1: C1--C4}

\begin{frame}{Slot 1 \textperiodcentered\ Thu 09:00--11:00}
  \begin{center}
    \begin{tabular}{@{}rl@{}}
      \toprule
      \textbf{Clock} & \textbf{Block} \\
      \midrule
      09:00--09:10 & Welcome, room rules, laptop check \\
      09:10--09:30 & \concept{C1} Agent vs chatbot \\
      09:30--09:55 & \concept{C2} Version control: your first pull request \\
      09:55--10:35 & \concept{C3} \artifact{AGENTS.md}: boundaries first, with privacy \\
      10:35--10:45 & Pause / set-up help \\
      10:45--11:00 & \concept{C4} To-do lists \\
      \bottomrule
    \end{tabular}
  \end{center}
\end{frame}

% ---------------------------------------------------------------- C1
\begin{frame}{C1 \textperiodcentered\ Agent vs chatbot}
  \timebox{09:10--09:30}
  \inplainwords{A chatbot only talks. An agent can also act: it opens your files,
    runs your code, reads the result and tries again.}
  \likethis{phoning a friend for advice vs hiring an RA who sits at your desk.}
\end{frame}

\begin{frame}{C1 \textperiodcentered\ Card--Krueger: same question, two answers}
  \centering
  {\large ``What is the NJ--PA change in employment?''}
  \vspace{0.6em}

  \begin{columns}[T,onlytextwidth]
    \begin{column}{0.47\textwidth}
      \begin{block}{Chatbot}
        Guesses a number from what it remembers about the 1994 paper.
      \end{block}
    \end{column}
    \begin{column}{0.47\textwidth}
      \begin{block}{Agent}
        Runs \artifact{did\_analysis.py} on our folder and quotes the printed note.
      \end{block}
    \end{column}
  \end{columns}
  \vspace{0.6em}

  Only one of these answers comes from \textbf{your} data.
\end{frame}

\begin{frame}{C1 \textperiodcentered\ Your turn}
  \youtry{read only, nothing is edited}
  \vspace{0.3em}
  \begin{enumerate}
    \item Ask a plain chat: ``What is the NJ--PA change in employment?''
    \item Ask the agent the same question, adding ``Do not edit any file.''
    \item Compare: where did each number come from?
  \end{enumerate}
  \vspace{0.5em}
  {\small Guide: \artifact{demos/slot1\_workflows.md} \textperiodcentered\ \artifact{exercises/04\_slot1\_workflows.md}}
\end{frame}

% ---------------------------------------------------------------- C2
\begin{frame}{C2 \textperiodcentered\ Version control}
  \timebox{09:30--09:55}
  \inplainwords{A record of every saved version of the project, each with a note,
    so you can see what changed and go back. GitHub keeps that record online and
    lets others review changes before they are accepted.}
  \likethis{Track Changes plus dated ``save as'' copies, shared with coauthors.}
\end{frame}

\begin{frame}{C2 \textperiodcentered\ Five words we use all day}
  \begin{tabular}{@{}p{0.27\linewidth}p{0.66\linewidth}@{}}
    \toprule
    \textbf{Word} & \textbf{Meaning} \\
    \midrule
    Project folder & All the course files, with their full history \\
    Saved version & A snapshot of the files, with a one-line note \\
    Separate copy & A safe place to try changes (a ``branch'') \\
    List of changes & The lines added and removed, side by side \\
    Pull request & A request to merge changes, with a review page \\
    \bottomrule
  \end{tabular}
\end{frame}

\begin{frame}{C2 \textperiodcentered\ Card--Krueger: reading the agent's edit}
  \begin{itemize}
    \item The agent edits \artifact{examples/card-krueger-toy/README.md}.
    \item GitHub shows a list of changes: green lines added, red lines removed.
    \item You accept, or ask for a fix, before it joins the main version.
    \item Nothing is lost: every saved version can come back.
  \end{itemize}
\end{frame}

\begin{frame}{C2 \textperiodcentered\ Your turn: first pull request}
  \youtry{09:30--09:55}
  \vspace{0.3em}
  \begin{enumerate}
    \item Make a separate copy.
    \item Edit one line in \artifact{notes/lane\_choice.md}: your tool and why.
    \item Save the version with a short note.
    \item Open a pull request and read its list of changes.
    \item Review, then merge.
  \end{enumerate}
  {\small Guide: \artifact{sessions/slot1.md} \textperiodcentered\ \artifact{exercises/05\_slot1\_governance.md}}
\end{frame}

% ---------------------------------------------------------------- C3
\begin{frame}{C3 \textperiodcentered\ \texttt{AGENTS.md}: boundaries first}
  \timebox{09:55--10:35}
  \inplainwords{A plain-text rules file the agent reads before it starts any task.
    It says what the agent may touch, what it must never do, and how to check its work.}
  \likethis{the house rules you give a new RA on day one.}
\end{frame}

\begin{frame}{C3 \textperiodcentered\ Card--Krueger: the agent refuses}
  Two lines from our \artifact{AGENTS.md}:
  \vspace{0.3em}
  \begin{block}{}
    \texttt{- Do not modify raw data.}\\
    \texttt{- Do not let agents certify identification, causality, or interpretation.}
  \end{block}
  \vspace{0.4em}
  Demo: ask the agent to edit \artifact{toy\_fast\_food.csv}.\\
  Watch it refuse and point to the rule.
\end{frame}

\begin{frame}{C3 \textperiodcentered\ Privacy in three questions}
  \begin{enumerate}
    \item How sensitive is the data?
    \item Who may see it?
    \item Where does the agent run: your laptop or someone else's computer?
  \end{enumerate}
  \vspace{0.5em}
  \begin{block}{Faculty default}
    \centering\large Anonymize \textperiodcentered\ authorize \textperiodcentered\ keep it on the machine
  \end{block}
  \vspace{0.3em}
  {\small Full check: \artifact{docs/privacy\_boundary.md}}
\end{frame}

\begin{frame}{C3 \textperiodcentered\ Your turn}
  \youtry{09:55--10:35}
  \vspace{0.3em}
  \begin{enumerate}
    \item Read \artifact{AGENTS.md} (two minutes).
    \item Add one rule you would want for your own project.
    \item Fill in the privacy check in \artifact{docs/privacy\_boundary.md}.
  \end{enumerate}
  \vspace{0.5em}
  {\small Guide: \artifact{demos/slot1\_governance.md} \textperiodcentered\ Pause 10:35--10:45}
\end{frame}

% ---------------------------------------------------------------- C4
\begin{frame}{C4 \textperiodcentered\ To-do lists}
  \timebox{10:45--11:00}
  \inplainwords{Before acting, the agent writes the steps it plans to take as a
    checklist. You edit the list, then it works through it and ticks items off.}
  \likethis{a shopping list you check before anyone drives to the store.}
\end{frame}

\begin{frame}{C4 \textperiodcentered\ Card--Krueger: plan only}
  \prompttext{\raggedright Plan only, do not edit. Task: examples/card-krueger-goals/goals/ck-robustness-checks.md}
  \vspace{0.6em}
  \begin{itemize}
    \item[$\square$] Load the panel
    \item[$\square$] Check that every store appears before and after
    \item[$\square$] Compute NJ and PA group means
    \item[$\square$] Drop one store at a time and recompute
    \item[$\square$] Write a short note
  \end{itemize}
\end{frame}

\begin{frame}{C4 \textperiodcentered\ Your turn}
  \youtry{10:45--11:00}
  \vspace{0.3em}
  \begin{enumerate}
    \item Ask for the to-do list only.
    \item Delete or reorder one step.
    \item Let the agent do step 1 only, then stop.
    \item Paste the list in \artifact{notes/modes\_playbook.md}.
  \end{enumerate}
  \vspace{0.5em}
  {\small Template: \artifact{templates/agent\_task\_contract.md}}
\end{frame}
